\documentclass[10pt,a4paper]{amsart}
\usepackage[margin=19mm]{geometry}
\usepackage{amsmath,amssymb,mathtools}
\usepackage[scr=rsfs,cal=euler]{mathalfa}
\usepackage{xfrac}
\usepackage[unicode,hidelinks,bookmarks=false]{hyperref}
\newcommand{\ie}{i.e.\ }
\newcommand{\cf}{cf.\ }
\newcommand{\resp}{respectively\ }
\newcommand{\Subsection}{\subsection}
\providecommand{\nobreakdash}{\nobreak\hbox{-}}
\setlength{\emergencystretch}{2em}
\multlinegap0pt
\usepackage{bm}
\usepackage{bbm}
\usepackage{dsfont}
\usepackage{xspace}
\usepackage{cancel}
\usepackage{mathtools}
\usepackage{cleveref}
\usepackage[shortlabels]{enumitem}
\usepackage{amsthm}
\makeatletter
\@ifundefined{theorem}{\newtheorem{theorem}{Theorem}}{}
\@ifundefined{remark}{\newtheorem{remark}[theorem]{Remark}}{}
\makeatother
\DeclareMathOperator{\pd}{pdim}
\DeclareMathOperator{\reg}{reg}
\def\P{{\mathcal P}}
\newcommand{\C}{\mathcal{C}}
\newcommand{\F}{\mathbb{F}}
\title[Neural codes via homological invariants of polarized neural ]{Neural codes via homological invariants of polarized neural ideals}
\author{Selvi Kara, Ellie Lew}
\date{Source: arXiv:2602.16993v1 (CC BY 4.0)}
\begin{document}
\maketitle

\noindent\emph{Source and licence.} Neural codes via homological invariants of polarized neural ideals. Version arXiv:2602.16993v1 (CC BY 4.0), distributed under the Creative Commons Attribution 4.0 International licence, as recorded in the arXiv licence metadata for this exact version and shown on the abstract page https://arxiv.org/abs/2602.16993v1. Full rights notice: licenses/arxiv-2602.16993-CC-BY-4.0.txt.\par\smallskip

\noindent\textit{Editorial scope.} Counted unit: the Theorem whose source label is \texttt{thm:reg1-classification} in the article (source0.tex lines 1595--1616, source label \texttt{thm:reg1-classification}), delivered together with the article's own complete original proof of it (source0.tex lines 1618--1697, one \texttt{proof} environment of 80 lines). No mathematics is written, repaired or re-derived: statement and proof are the article's own text line for line. Required context retained uncounted: rem:basic\_degree\_bounds (lines 710--726). Every definition, display and label of those lines is retained unchanged. Omitted from this package and not redistributed: 1--709, 727--1594, 1617--1617, 1698--2222. Nothing inside a retained excerpt is omitted or altered: every retained body is delivered exactly as the article's lines are written, so no omission receipt of any kind accompanies this package. Citations: the retained excerpts carry no citation command at all, so no bibliography entry of the article is transcribed and no reference is redistributed. Reference adaptation: none was needed and none was made; every reference inside the delivered unit resolves inside it, so no unresolved reference or citation is produced. Printed slips kept literal: no printed word, symbol, hypothesis or number of the counted statement or of its proof was altered. Layout and class: the article's own class is \texttt{amsart} with packages including geometry, xcolor, inputenc, amsfonts, amsthm, amssymb, amscd, amsmath, blindtext, caption, subcaption, pdflscape, multirow, graphicx; the wrapper is the lane's pinned \texttt{amsart} editorial wrapper at 10pt on A4, which restates the theorem-like environments the retained text needs and adds the article's own macro definitions that the retained text needs (\texttt{\textbackslash C}, \texttt{\textbackslash F}, \texttt{\textbackslash P}, \texttt{\textbackslash pd}, \texttt{\textbackslash reg}), with extra wrapper packages (bm, bbm, dsfont, xspace, cancel, mathtools, cleveref, enumitem). The delivered numbering is the unit's own. Assets: no retained span contains a figure environment, a graphics-inclusion command, a PDF-page inclusion, a table or a TikZ picture; no image file is copied into this package, the package declares no asset dependency and no image file is redistributed with it. Licence: version arXiv:2602.16993v1 of this article is distributed under the Creative Commons Attribution 4.0 International licence, as recorded in the arXiv licence metadata for this exact version and shown on the abstract page https://arxiv.org/abs/2602.16993v1. A full rights notice is delivered at licenses/arxiv-2602.16993-CC-BY-4.0.txt. Characters: every retained excerpt is pure ASCII, so the delivered engine, XeLaTeX with its default Latin Modern text font, reports no missing character.
\begin{remark}\label{rem:basic_degree_bounds}
Let $\C\subseteq\F_2^n$ with polarized neural ideal $I=\P(J_\C)$ and polar complex $\Delta=\Delta_\C$.
\begin{enumerate}[label=(\alph*)]
\item No minimal generator of $I$ is divisible by $x_i y_i$.

\item If $x_i,y_i\in I$, then depolarization gives $x_i, 1-x_i\in J_\C$. Hence $1\in J_\C$.
Thus $J_\C=R$ and $\C=V(J_\C)=\emptyset$.

\item The ideal $I$ is squarefree and is supported on a vertex set
$V(\Delta)\subseteq\{x_1,\dots,x_n,y_1,\dots,y_n\}$, so $|V(\Delta)|\le 2n$.
In particular, $\beta_{i,j}(I)=0$ for all $j>2n$.

\item Every minimal generator of $I$ is a \emph{transversal} monomial $x_\sigma y_\tau$ with
$\sigma\cap\tau=\varnothing$.  Hence $\deg(x_\sigma y_\tau)=|\sigma|+|\tau|\le n$, and therefore
$\beta_{0,j}(I)=0$ for all $j>n$.
\end{enumerate}
\end{remark}

\begin{theorem}\label{thm:reg1-classification}
Let $\C\subset \F_2^n$ be a nonempty neural code.
Then the following are equivalent:
\begin{enumerate}
\item $\reg(I)=1$;
\item $I$ is generated by variables and, after relabeling,
\[
I=\langle x_i : i\in\sigma\rangle+\langle y_j : j\in\tau\rangle
\qquad\text{for some disjoint }\sigma,\tau\subseteq[n];
\]
\item $\C$ is a subcube of $\F_2^n$ obtained by fixing some coordinates to $0$ and some disjoint coordinates to $1$, i.e.
\[
\C=\{ v\in \F_2^n : v_i=0\ \forall i\in\sigma,\ \ v_j=1\ \forall j\in\tau \}
\]
for some disjoint $\sigma,\tau\subseteq[n]$.
\end{enumerate}
Moreover, if these conditions hold and $t:=|\sigma|+|\tau|$, then
\[
\pd(I)=t-1\le n-1,
\]
and $\pd(I)=n-1$ if and only if $t=n$, equivalently if and only if $\C$ consists of a single codeword.
\end{theorem}

\begin{proof}
\emph{$(1)\Rightarrow(2)$:}
If $\reg(I)=1$, then every minimal monomial generator of $I$ has degree $1$. Hence $I$ is generated by variables.
Since $I$ is proper, it cannot contain both $x_i$ and $y_i$ for any $i$ by \Cref{rem:basic_degree_bounds}(2).
Therefore
\[
I=\langle x_i:i\in\sigma\rangle+\langle y_j:j\in\tau\rangle
\]
for some disjoint $\sigma,\tau\subseteq[n]$.


\emph{$(2)\Rightarrow(1)$:}
Conversely, any ideal generated by variables has a $1$--linear resolution and $\reg(I)=1$.


\emph{$(2)\Rightarrow(3)$:}
Assume $I=\langle x_i:i\in\sigma\rangle+\langle y_j:j\in\tau\rangle$ with $\sigma\cap\tau=\varnothing$, and set
\[
Q:=\{ v\in\F_2^n : v_i=0\ \forall i\in\sigma,\ v_j=1\ \forall j\in\tau \}.
\]
Since $x_i\in I$, then $x_i\in J_\C$. So $x_i(c)=c_i=0$ for all $c\in\C$.
Similarly, $y_j\in I$ implies that $(1-x_j)\in J_\C$. So $(1-x_j)(c)=1-c_j=0$ for all $c\in\C$.
Thus $\C\subseteq Q$.

To prove equality, suppose $\C\subsetneq Q$ and choose $u\in Q\setminus \C$.
Then $u$ is a noncodeword, so $\rho_u\in J_\C$.
Choose a minimal pseudo-monomial $f\in J_\C$ dividing $\rho_u$. By definition
$f\in CF(J_\C)$.

We claim that $\deg(f)\ge 2$.
Indeed, since $u\in Q$, for each $i\in\sigma$ one has $u_i=0$, so $\rho_u$ is divisible by $(1-x_i)$ but
not by $x_i$, and for each $j\in\tau$ one has $u_j=1$, so $\rho_u$ is divisible by $x_j$ but not by $(1-x_j)$.
Thus none of the linear pseudo-monomials $x_i$ ($i\in\sigma$) or $(1-x_j)$ ($j\in\tau$) divides $\rho_u$.
If $k\notin\sigma\cup\tau$, then neither $x_k$ nor $(1-x_k)$ lies in $J_\C$, since otherwise the $k$th
coordinate would be fixed on $Q$, contradicting the description of $Q$.
Hence no degree-$1$ pseudo-monomial in $J_\C$ divides $\rho_u$. Therefore $\deg(f)\ge 2$.

Notice that $\P(f)\in I$.
On the other hand, none of the variable generators of $I$ in~(2) divides $\P(f)$: the same divisibility
check as above shows that neither $x_i$ ($i\in\sigma$) nor $y_j$ ($j\in\tau$) divides $\P(f)$.
This contradicts the assumption that $I$ is generated by $\{x_i:i\in\sigma\}\cup\{y_j:j\in\tau\}$.
Therefore $\C=Q$.


\emph{$(3)\Rightarrow(2)$:}
Assume $\C$ as in~(3), with $\sigma$ and $\tau$ the sets of coordinates fixed to $0$ and $1$.
Every noncodeword violates at least one fixed coordinate, so each characteristic pseudo-monomial $\rho_v$ is divisible by
some $x_i$ ($i\in\sigma$) or some $(1-x_j)$ ($j\in\tau$). Hence
\[
J_\C\subseteq \langle x_i:i\in\sigma,\ (1-x_j):j\in\tau\rangle.
\]
Conversely, for $i\in\sigma$ we have
\[
\sum_{v:  v_i=1}\rho_v
 = 
x_i\prod_{k\neq i}\bigl(x_k+(1-x_k)\bigr)
 = x_i\in J_\C,
\]
and for $j\in\tau$,
\[
\sum_{v:  v_j=0}\rho_v
 = 
(1-x_j)\prod_{k\neq j}\bigl(x_k+(1-x_k)\bigr)
 = 1-x_j\in J_\C.
\]
Thus
\[
J_\C=\langle x_i:i\in\sigma \rangle +\langle (1-x_j):j\in\tau\rangle,
\]
and polarizing gives
\[
I=\langle x_i:i\in\sigma\rangle+\langle y_j:j\in\tau\rangle.
\]


Finally, if $I$ is generated by $t=|\sigma|+|\tau|$ variables, then the Taylor on these variables
yields a  minimal free resolution of $I$ of
length $t-1$. Therefore $\pd(I)=t-1\le n-1$. Moreover, $t=n$ if and only if every coordinate is fixed,
equivalently if and only if $\C$ consists of a single codeword.
\end{proof}

\end{document}
